% 作者 shingja · github.com/shingja/shingja-math-modeling-workbench · 双许可见 LICENSE · 含隐形溯源标记
% !TEX program = xelatex
% 2026 CUMCM 数学建模论文 LaTeX 模板
% Aligned with: 全国大学生数学建模竞赛论文格式规范(2026版)

\documentclass[12pt,a4paper]{ctexart}

% === Packages ===
\usepackage[top=2.5cm, bottom=2.5cm, left=2.5cm, right=2.5cm]{geometry}
\usepackage{amsmath, amssymb, amsthm}
\usepackage{graphicx}
\usepackage{booktabs}       % 三线表
\usepackage{hyperref}
\usepackage{listings}       % 代码附录
\usepackage{xcolor}
\usepackage{fancyhdr}
\usepackage{titling}
\usepackage{caption}
\usepackage{float}

% === Code listing style ===
\lstset{
    basicstyle=\small\ttfamily,
    breaklines=true,
    frame=single,
    numbers=left,
    numberstyle=\tiny,
    commentstyle=\color{gray},
    keywordstyle=\color{blue},
    stringstyle=\color{red},
}

% === Page style: plain = page number centered at footer, no header ===
\pagestyle{plain}
\pagenumbering{arabic}

% === Macros ===
% Three-line table
\newcommand{\threetab}[1]{%
    \begin{table}[H]
    \centering
    \caption{#1}
    \begin{tabular}
}
\newcommand{\endthreetab}{%
    \end{tabular}
    \end{table}
}

% Publication-quality figure
\newcommand{\pubfig}[3]{%
    \begin{figure}[H]
    \centering
    \includegraphics[width=\textwidth, keepaspectratio]{#1}
    \caption{#2}
    \label{#3}
    \end{figure}
}

% --- DOCUMENT START ---
\begin{document}

% ================ 第 1 页: 承诺书 ================
\thispagestyle{empty}
% 作者 shingja · github.com/shingja/shingja-math-modeling-workbench · 双许可见 LICENSE · 保留版权标识
% 承诺书 (2026 CUMCM 固定内容, 不编号)
\begin{center}
    {\LARGE \textbf{承\hspace{1em}诺\hspace{1em}书}}
\end{center}

\vspace{2cm}

我们仔细阅读了《全国大学生数学建模竞赛章程》和《全国大学生数学建模竞赛论文格式规范》(2026版)。

我们完全明白，在竞赛开始后参赛队员不能以任何方式（包括电话、电子邮件、网上咨询等）与队外的任何人（包括指导教师）研究、讨论与赛题有关的问题。

我们知道，抄袭别人的成果是违反竞赛章程和参赛规则的，如果引用别人的成果或其他公开的资料（包括网上查到的资料），必须按照规定的参考文献的表述方式在正文引用处和参考文献中明确列出。

我们郑重承诺，严格遵守竞赛章程和参赛规则，以保证竞赛的公正、公平性。如有违反竞赛章程和参赛规则的行为，我们将受到严肃处理。

\vspace{2cm}

我们参赛选择的题号是（从A/B/C/D/E中选择一项填写）：\underline{\hspace{3cm}}

\vspace{1cm}

\begin{flushright}
    参赛队员（签名）：\underline{\hspace{2cm}} \\
    \vspace{0.5cm}
    日期：\underline{\hspace{2cm}}年\underline{\hspace{1cm}}月\underline{\hspace{1cm}}日
\end{flushright}

\newpage

% ================ 第 2 页: 编号专用页 ================
\thispagestyle{empty}
% 作者 shingja · github.com/shingja/shingja-math-modeling-workbench · 双许可见 LICENSE · 保留版权标识
% 编号专用页 (2026 CUMCM)
\begin{center}
    {\LARGE \textbf{赛\hspace{1em}区\hspace{1em}评\hspace{1em}阅\hspace{1em}编\hspace{1em}号}}
\end{center}

\vspace{3cm}

\begin{center}
    \begin{tabular}{|c|c|c|c|c|c|}
    \hline
    & & & & & \\
    \hline
    & & & & & \\
    \hline
    & & & & & \\
    \hline
    \end{tabular}
\end{center}

\vspace{2cm}

\begin{flushright}
    全国大学生数学建模竞赛组委会\\
    2026年9月
\end{flushright}

\newpage

% ================ 第 3 页: 摘要专用页 ================
\setcounter{page}{1}  % 页码从 1 开始

\begin{center}
    {\LARGE \textbf{论文标题}}
\end{center}

\vspace{1cm}

\begin{abstract}
\noindent
\textbf{摘要：}本文针对[问题背景]，建立了[模型名称]模型。\\
针对问题一，采用[方法]，求解得到[具体数值结果]；\\
针对问题二，采用[方法]，求解得到[具体数值结果]；\\
\ldots

结果表明，[模型]能够[结论]。通过灵敏度分析[验证方式]验证了模型的[稳定性/准确性]。
\end{abstract}

\vspace{0.5cm}
\noindent\textbf{关键词：}关键词1；关键词2；关键词3；关键词4
\newpage

% ================ 正文 ================
% 作者 shingja · github.com/shingja/shingja-math-modeling-workbench · 双许可见 LICENSE · 保留版权标识
% ============================================================
% 01 — 问题重述
% 说明: 用自己的语言重新组织问题背景, 不得照抄原文。
%       明确列出待解决的各个子问题。
%       数值必须从 results/results.json 提取, 禁止手写。
% ============================================================

\section{问题重述}

\subsection{问题背景}
% [此处填写问题背景。用自己的语言概括赛题给出的实际场景和研究动机。]

本题来源于[实际应用领域]中的[核心问题]。在[具体场景]中，需要[研究目标]。该问题的解决对于[实际意义]具有重要意义。

\subsection{问题提出}
% [此处列出需要解决的子问题, 条理清晰。]

基于上述背景，本文需要解决以下问题：

\textbf{问题一：}[问题一简要描述]

\textbf{问题二：}[问题二简要描述]

\textbf{问题三：}[问题三简要描述]

\textbf{问题四：}[问题四简要描述]

% 作者 shingja · github.com/shingja/shingja-math-modeling-workbench · 双许可见 LICENSE · 保留版权标识
% ============================================================
% 02 — 模型假设
% 说明: 列出建模所做的合理简化假设, 每条编号。
%       假设应具体、可验证, 并简要说明理由。
%       数值必须从 results/results.json 提取, 禁止手写。
% ============================================================

\section{模型假设}

为简化问题并建立可求解的数学模型，本文做出以下合理假设：

\begin{enumerate}
    \item \textbf{[假设名称]:} [假设内容及合理性说明]
    \item \textbf{[假设名称]:} [假设内容及合理性说明]
    \item \textbf{[假设名称]:} [假设内容及合理性说明]
    \item \textbf{[假设名称]:} [假设内容及合理性说明]
    \item \textbf{[假设名称]:} [假设内容及合理性说明]
\end{enumerate}

以上假设在简化模型的同时保留了问题的核心特征，后续章节将在此基础上展开建模。

% 作者 shingja · github.com/shingja/shingja-math-modeling-workbench · 双许可见 LICENSE · 含隐形溯源标记
% ============================================================
% 03 — 符号说明
% 说明: 使用三线表统一列出文中使用的数学符号。
%       符号按出现顺序排列, 注明含义和单位。
%       数值必须从 results/results.json 提取, 禁止手写。
% ============================================================

\section{符号说明}

本文使用的主要数学符号及其含义如表\ref{tab:symbols}所示。

\begin{table}[H]
    \centering
    \caption{主要符号说明}
    \label{tab:symbols}
    \begin{tabular}{ccl}
        \toprule
        \textbf{符号} & \textbf{单位} & \textbf{含义} \\
        \midrule
        $x$      & —       & [符号含义] \\
        $y$      & —       & [符号含义] \\
        $t$      & s       & [符号含义] \\
        $\alpha$ & —       & [符号含义] \\
        $\beta$  & —       & [符号含义] \\
        \bottomrule
    \end{tabular}
\end{table}

注：未在表中列出的符号将在首次出现时予以说明。

% 作者 shingja · github.com/shingja/shingja-math-modeling-workbench · 双许可见 LICENSE · 保留版权标识
% ============================================================
% 04 — 问题一: 模型建立与求解
% 说明: 按照 "建模思路 → 模型建立 → 算法设计 → 求解结果" 结构。
%       关键数值必须从 results/results.json 提取, 禁止手写。
%       数值格式: 4位有效数字, 百分比2位小数, p值3位小数。
% ============================================================

\section{问题一: [问题一标题]}

\subsection{建模思路}
% [此处描述问题分析过程和建模思路, 说明为什么选择该方法。]

针对问题一，首先分析[核心矛盾/关键特征]。考虑到[约束条件/数据特点]，本文采用[方法名称]进行建模，主要基于以下考虑：(1) [理由一]; (2) [理由二]; (3) [理由三]。

\subsection{模型建立}
% [此处给出数学模型: 目标函数/约束条件/变量定义。使用 align 环境排版关键公式。]

定义决策变量 $x = (x_1, x_2, \ldots, x_n)^T$，建立如下数学模型：

\begin{equation}
    \min_{x} \quad f(x) = \text{[目标函数表达式]}
\end{equation}

\begin{equation}
    \text{s.t.} \quad
    \begin{cases}
        g_i(x) \leq 0, \quad i = 1, 2, \ldots, m \\
        h_j(x) = 0, \quad j = 1, 2, \ldots, p \\
        x \in \Omega
    \end{cases}
\end{equation}

其中，$g_i(x)$ 为不等式约束，$h_j(x)$ 为等式约束，$\Omega$ 为决策变量的可行域。

\subsection{算法设计与求解}
% [此处描述求解算法: 算法流程/参数设置/收敛判据。]

本文采用[算法名称]求解上述模型，算法流程如下：

\begin{enumerate}
    \item \textbf{初始化:} [参数初始化说明]
    \item \textbf{迭代求解:} [迭代过程说明]
    \item \textbf{终止判据:} [收敛条件说明]
\end{enumerate}

\subsection{结果分析}
% [此处用文字解读结果, 引用表格和图形。数值从 results.json 提取。]

求解得到[关键结果]。如表\ref{tab:prob1_result}和图\ref{fig:prob1}所示，[结果解读]。

\begin{table}[H]
    \centering
    \caption{问题一主要结果}
    \label{tab:prob1_result}
    \begin{tabular}{ccc}
        \toprule
        \textbf{指标} & \textbf{数值} & \textbf{单位} \\
        \midrule
        [指标1] & [数值] & [单位] \\
        [指标2] & [数值] & [单位] \\
        \bottomrule
    \end{tabular}
\end{table}

% 作者 shingja · github.com/shingja/shingja-math-modeling-workbench · 双许可见 LICENSE · 保留版权标识
% ============================================================
% 05 — 问题二: 模型建立与求解
% 说明: 结构与问题一类似, 注意与问题一的递进关系。
%       若问题二在问题一基础上扩展, 需明确指出继承与改进之处。
%       数值必须从 results/results.json 提取, 禁止手写。
%       数值格式: 4位有效数字, 百分比2位小数, p值3位小数。
% ============================================================

\section{问题二: [问题二标题]}

\subsection{建模思路}
% [此处描述与问题一的关联和问题二的特殊性。]

问题二在问题一的基础上，进一步考虑了[新增因素/放宽假设]。这使得原模型在[某方面]不再适用，需要对模型进行[扩展/修正]。本文采用[方法名称]来解决这一问题。

\subsection{模型建立}
% [此处给出扩展后的数学模型。]

在问题一模型的基础上，引入[新变量/新约束]，建立如下扩展模型：

\begin{equation}
    \min_{x, y} \quad F(x, y) = \text{[目标函数表达式]}
\end{equation}

\begin{align}
    \text{s.t.} \quad
    & g_i(x, y) \leq 0, \quad i = 1, 2, \ldots, m' \\
    & h_j(x, y) = 0, \quad j = 1, 2, \ldots, p' \\
    & x \in \Omega_x, \; y \in \Omega_y
\end{align}

\subsection{算法设计与求解}
% [此处描述求解算法, 说明与问题一算法的异同。]

考虑到问题二[计算复杂度/特殊结构]，本文采用[算法名称]进行求解。相较于问题一，算法在[某方面]做了[改进/调整]。

\subsection{结果分析}
% [此处解读结果, 并与问题一结果进行对比分析。]

求解得到[关键结果]。与问题一结果对比（见表\ref{tab:comparison}），[对比分析结论]。

\begin{table}[H]
    \centering
    \caption{问题一与问题二结果对比}
    \label{tab:comparison}
    \begin{tabular}{cccc}
        \toprule
        \textbf{指标} & \textbf{问题一} & \textbf{问题二} & \textbf{变化} \\
        \midrule
        [指标1] & [数值] & [数值] & [变化量/率] \\
        [指标2] & [数值] & [数值] & [变化量/率] \\
        \bottomrule
    \end{tabular}
\end{table}

% 作者 shingja · github.com/shingja/shingja-math-modeling-workbench · 双许可见 LICENSE · 保留版权标识
% ============================================================
% 06 — 灵敏度分析
% 说明: 检验关键参数变化对模型结果的影响程度。
%       至少包含参数扰动分析和鲁棒性检验。
%       数值必须从 results/results.json 提取, 禁止手写。
%       数值格式: 4位有效数字, 百分比2位小数, p值3位小数。
% ============================================================

\section{灵敏度分析}

为验证模型的稳定性和鲁棒性，本节对模型中的关键参数进行灵敏度分析。

\subsection{参数扰动分析}
% [此处分析关键参数在合理范围内变化时, 目标函数/决策变量的变化。]

选取模型中[参数1]、[参数2]和[参数3]作为关键参数，在基准值的 $\pm 10\%$ 范围内进行扰动分析。结果如表\ref{tab:sensitivity}所示。

\begin{table}[H]
    \centering
    \caption{关键参数灵敏度分析}
    \label{tab:sensitivity}
    \begin{tabular}{ccccc}
        \toprule
        \textbf{参数} & \textbf{变化范围} & \textbf{目标值变化} & \textbf{变化率} & \textbf{灵敏度等级} \\
        \midrule
        [参数1] & $[-10\%, +10\%]$ & [范围] & [百分比] & [低/中/高] \\
        [参数2] & $[-10\%, +10\%]$ & [范围] & [百分比] & [低/中/高] \\
        [参数3] & $[-10\%, +10\%]$ & [范围] & [百分比] & [低/中/高] \\
        \bottomrule
    \end{tabular}
\end{table}

分析结果表明，[参数1]对模型结果影响[程度]，灵敏度等级为[低/中/高]；而[参数2]和[参数3]的影响[程度]。

\subsection{鲁棒性检验}
% [此处从不同角度验证模型的可靠性: 不同初始值/不同数据子集/不同随机种子等。]

为进一步验证模型的鲁棒性，本文进行了以下检验：

\begin{enumerate}
    \item \textbf{初始值敏感性:} 使用[数量]组不同初始值重复求解，目标函数值的标准差为[数值]（相对标准差为[百分比]），表明算法对初始值[不/较/非常]敏感。
    \item \textbf{数据扰动检验:} 对输入数据添加[水平]的高斯噪声，重复[次数]次实验，结果变化在[范围]内，表明模型对数据噪声[鲁棒性程度]。
\end{enumerate}

综合灵敏度分析结果，模型在关键参数扰动下表现[稳定/较稳定/不稳定]，验证了模型的[可靠性/有效性]。

% 作者 shingja · github.com/shingja/shingja-math-modeling-workbench · 双许可见 LICENSE · 保留版权标识
% ============================================================
% 07 — 模型评价
% 说明: 客观评价模型优缺点, 提出改进方向。
%       分为"模型优点""模型缺点""模型改进与推广"三部分。
%       数值必须从 results/results.json 提取, 禁止手写。
% ============================================================

\section{模型评价与改进}

\subsection{模型优点}

本文建立的数学模型具有以下优点：

\begin{enumerate}
    \item \textbf{[优点1标题]:} [具体说明, 结合模型特点和数据结果]
    \item \textbf{[优点2标题]:} [具体说明, 结合模型特点和数据结果]
    \item \textbf{[优点3标题]:} [具体说明, 结合模型特点和数据结果]
    \item \textbf{[优点4标题]:} [具体说明, 结合模型特点和数据结果]
\end{enumerate}

\subsection{模型缺点}

本模型仍存在以下不足：

\begin{enumerate}
    \item \textbf{[缺点1标题]:} [具体说明, 客观分析]
    \item \textbf{[缺点2标题]:} [具体说明, 客观分析]
    \item \textbf{[缺点3标题]:} [具体说明, 客观分析]
\end{enumerate}

\subsection{模型改进与推广}

针对上述不足，未来可从以下方向进行改进：

\begin{enumerate}
    \item \textbf{算法改进:} [改进思路, 如采用更高效的算法/并行计算等]
    \item \textbf{模型扩展:} [扩展思路, 如考虑更多实际因素/放宽假设等]
    \item \textbf{应用推广:} [推广方向, 如将模型应用于其他领域/场景]
\end{enumerate}


% ================ 参考文献 ================
\begin{thebibliography}{99}

\bibitem{ref1} 作者. 题名[J]. 期刊, 年份, 卷(期): 页码. DOI:xxx

\end{thebibliography}

% ================ 附录: 全部源代码 ================
\newpage
\appendix
\section{程序源代码}

\subsection{问题一求解程序}
\lstinputlisting[language=Matlab]{../../matlab/templates/skeleton.m}

\end{document}
